\scanpage{128}
\begin{center}- \mspage{118} -\end{center}\margin{IX}

de type fini. On constate alors que l'homomorphisme (9.5.6) est l'homomorphisme canonique qui exprime l'action du sous-groupe d'inertie $I\subset\pi$ sur $X_\eta(\bar K)$, défini par la méthode de 9.2. Pour le vérifier, on se ramène immédiatement à la situation de 9.4, en y supposant que $n$ est égal à une puissance de $\ell$, et que $R$ est un $\Lambda$-Module localement constant (tout comme $Q$ et $P$), de sorte qu'il en est également de même de $E$, $F$ et de $X_\eta$. La vérification de l'identité des deux définitions est alors laissée au lecteur.

9.6. Nous sommes maintenant en mesure de définir l'accouplement canonique annoncé (9.1.2), sans restriction sur $\ell$, de telle façon que pour $\ell\neq p$ cette définition coïncide avec celle donnée dans 9.2. Pour ceci, on applique les considérations de 9.5 à la situation suivante :
\[
(9.6.1)\quad\left\{\begin{array}{l}
P=T_\ell(A^0)^t=\check{\underline M}'_\ell(1),\quad Q=M_\ell,\quad R=T_\ell(A^0)^f/T_\ell(A^0)^t,\\[2em]
F=T_\ell(A^0)^f,\quad E=D(F')=D(T_\ell(A'^0)^f),
\end{array}\right.
\]
où dans la dernière formule $D$ désigne le dual de Cartier au sens des pro-$\ell$-groupes de Barsotti-Tate, i.e $D(Y)$ est le système \uline{projectif} des $D(Y(n))$, pour les morphismes de transition habituels. La structure de $F$ comme extension de $R$ par $P$ est claire. On a de même sur $F'=T_\ell(A'^0)^f$ une structure d'extension de $R'$ par $P'=T_\ell(A'^0)^t=\check{\underline M}_\ell(1)$. Appliquant la dualité de Cartier, on en conclut que $E=D(F')$ est muni d'une structure d'extension de $D(P')=M_\ell$ par $D(R')$. Or $D(R')$ est canoniquement isomorphe à $R$ en vertu de 5.5 (prouvé dans 7.4.6), ce qui précise la structure de $E$ comme extension de $Q$ par $R$.

\margin{\unsure{436|426}}\margin{436}\margin{127}
